\documentclass[10pt]{article}
\usepackage[utf8]{inputenc}
\usepackage[T1]{fontenc}
\usepackage{amsmath}
\usepackage{amsfonts}
\usepackage{amssymb}
\usepackage[version=4]{mhchem}
\usepackage{stmaryrd}
\usepackage{bbold}
\usepackage{hyperref}
\hypersetup{colorlinks=true, linkcolor=blue, filecolor=magenta, urlcolor=cyan,}
\urlstyle{same}

\begin{document}
\section*{Job market signaling.}
Spence (1973), Spence (1974), Hörner (2008), Mas-Colell et al. (1995).

\section*{Outline of the game}
\begin{itemize}
  \item Nature determines whether an agent has low or high productivity.
  \item The agent privately observes his/her productivity and choose a level of education. (Player 1)
  \item The "labor market" (i.e. player 2) observes the education chosen by the worker and offers a wage.
  \item Worker accept or reject the offer.
\end{itemize}

\section*{Spence's job market signaling}
\begin{itemize}
  \item Nature chooses the agent's productivity $\eta \in\left\{\eta_{\ell}, \eta_{h}\right\}, 0<\eta_{\ell}<\eta_{h}$; with probability $\lambda \in(0,1), \eta=\eta_{h}$.
  \item After privately observing $\eta$, the agent chooses $e \in\{0,1\}$. Education, that is $e=1$, costs $c_{\ell}$ to type $\eta_{\ell}$ and $c_{h}$ to type $\eta_{h}$; no education ( $e=0$ ) costs zero to both types.
  \item After observing $e$, the market or the employer offers a wage $w \in \mathbb{R}$ and the agent accepts or rejects the offer.
  \item A worker that receives $w$ and pays $c$ has a payoff $w-c$.
\end{itemize}

If not employed, the agent's payoff is $0-c$.

\begin{itemize}
  \item An employer who hires a worker with productivity $\eta$ and pays $w$ has payoff $\eta-w$.
\end{itemize}

If no worker is hired, the employer's payoff is zero. Employers are risk neutral.

A signaling game

\section*{Player 2}
Various alternatives to model ˋˋthe market."

\begin{itemize}
  \item Reduced form. The market wage will satisfy a zero expected profit condition where expectation is taken with respect to the probability generated by $\lambda, \sigma_{1}$.
  \item Explicit. After observing $e$, two (or more) firms simultaneously and independently offer a wage. Bertrand competition then insures a zero expected profit condition in any sequential equilibrium.
  \item Fictitious player 2 with payoff $U_{2}(\eta, e, w)=-(\eta-w)^{2}$
\end{itemize}

\section*{Benchmark}
\begin{itemize}
  \item No information. No player observes $\eta$.
\end{itemize}

Equilibrium: $\left(e=0, w=(1-\lambda) \eta_{\ell}+\lambda \eta-h\right)$.

\begin{itemize}
  \item Complete information. Both players observe $\eta$.
\end{itemize}

Equilibrium: $(e=0, w=\eta)$\\
Efficiency ex ante and efficiency ex post.

\section*{Separating equilibium}
Proposed equilibrium.

\begin{itemize}
  \item Player 1. $\left(s_{1}\left(\eta_{l}\right), s_{1}\left(\eta_{h}\right)\right)=(0,1)$
  \item Player 2. $\left(s_{2}(0), s_{2}(1)\right)=\left(w_{0}, w_{1}\right)=\left(\eta_{\ell}, \eta_{h}\right)$\\
$s_{2}$ is sequentially rational.
  \item Player 2's information sets $e=0$ and $e=1$ are reached with probability $(1-\lambda)>0, \lambda>0$ respectively.
  \item Player 2's beliefs are $\mu\left(\eta_{\ell} \mid e=0\right)=1, \mu\left(\eta_{h} \mid e=1\right)=1$. Bayes rule.
  \item Zero profit implies $w_{0}=\eta_{\ell}, w_{1}=\eta_{h}$.
\end{itemize}

It remains to show that $s_{1}$ is sequentially rational

\section*{Incentive compatibility}
It remains to show that $s_{1}$ is sequentially rational

\begin{itemize}
  \item player 1 , type $\eta_{\ell}$.
\end{itemize}

$$
\begin{aligned}
U_{1}\left(\eta_{\ell}, e=0, s_{2}(0)\right) & \geq U_{1}\left(\eta_{\ell}, e=1, s_{2}(1)\right) \\
\eta_{\ell} & \geq \eta_{h}-c_{\ell}
\end{aligned}
$$

Player $1, \eta_{\ell}$ incentives are compatible with proposed strategies if the inequality is satisfied.

\begin{itemize}
  \item player 1, type $\eta_{h}$.
\end{itemize}

$$
\begin{aligned}
U_{1}\left(\eta_{h}, e=1, s_{2}(1)\right) & \geq U_{1}\left(\eta_{h}, e=0, s_{2}(0)\right) \\
\eta_{h}-c_{h} & \geq \eta_{\ell}
\end{aligned}
$$

Player $1, \eta_{h}$ with proposed strategies if the inequality is satisfied.\\
The proposed strategies are a separating equilibrium if

$$
\eta_{h}-c_{h} \geq \eta_{\ell} \geq \eta_{h}-c_{\ell}
$$

This implies $c_{h} \leq c_{\ell}$.

\section*{Pooling equilibium}
Proposed equilibrium

\begin{itemize}
  \item Player 1. $\left(s_{1}\left(\eta_{\ell}\right), s_{1}\left(\eta_{h}\right)\right)=(1,1)$
  \item Player 2. $\left(s_{2}(0), s_{2}(1)\right)=\left(w_{0}, w_{1}\right)=\left(\eta_{\ell}, \bar{\eta}_{b}\right)$ where
\end{itemize}

$$
\bar{\eta}_{b}=(1-\lambda) \eta_{\ell}+\lambda \eta_{h} .
$$

$s_{2}$ is sequentially rational.

\begin{itemize}
  \item Player 2's information sets $e=0$ is not on the path of play, $e=1$ is reached with probability 1 .
  \item Player 2's beliefs when $e=1$ are $\mu\left(\eta_{\ell} \mid e=1\right)=1-\lambda, \mu\left(\eta_{h} \mid e=1\right) =\lambda$. Bayes rule.
  \item Player 2's beliefs when $e=0$ are $\mu\left(\eta_{\ell} \mid e=0\right)=1$.
  \item Zero expected profit condition implies $w_{0}=\eta_{\ell}, w_{1}=\bar{\eta}_{b}$.
\end{itemize}

It remains to show that $s_{1}$ is sequentially rational

\section*{Pooling equilibrium}
It remains to show that $s_{1}$ is sequentially rational.\\
Player 1, type $\eta_{\ell}$ :

$$
\begin{aligned}
U_{1}\left(\eta_{\ell}, e=1, s_{2}(1)\right) & \geq U_{1}\left(\eta_{\ell}, e=0, s_{2}(0)\right) \\
\bar{\eta}_{b}-c_{\ell} & \geq \eta_{\ell}
\end{aligned}
$$

Player 1, type $\eta_{h}$ :

$$
\begin{aligned}
U_{1}\left(\eta_{h}, e=1, s_{2}(1)\right) & \geq U_{1}\left(\eta_{h}, e=0, s_{2}(0)\right) \\
\bar{\eta}_{b}-c_{h} & \geq \eta_{\ell}
\end{aligned}
$$

\section*{Remarks}
\begin{itemize}
  \item Separating equilibrium is suboptimal. Education is not an argument in payoff functions.
  \item P.O. allocation is pooling equilibrium without education.
  \item At P.O. allocation high-productivity types subsidize low productivity types.
\end{itemize}

\section*{References}
Hörner, Johannes. 2008. "Signalling and Screening." In New Palgrave Dictionary of Economics Living Edition, edited by Matias Vernengo and Barkley J. Rosser Jr Esteban Perez Caldentey and. Springer. \href{https://doi.org/doi.org/10.1057/978-1-349-95121-5_2097-1}{https://doi.org/doi.org/10.1057/978-1-349-95121-5\_2097-1}.\\
Mas-Colell, Andreu, Michael Whinston, and Jerry Green. 1995. Microeconomic Theory. Oxford University Press.\\
Spence, A. Michael. 1973. "Job Market Signaling." The Quarterly Journal of Economics 87 (3): 355-74. \href{https://doi.org/10.2307/1882010}{https://doi.org/10.2307/1882010}.\\
Spence, A. Michael. 1974. Market Signaling. Harvard University Press.


\end{document}